\documentclass[11pt,a4paper]{article}
% Preâmbulo comum. Não duplicar pacotes nos módulos.
\usepackage[utf8]{inputenc}
\usepackage[T1]{fontenc}
\usepackage[brazil]{babel}
\usepackage{amsmath,amssymb,amsthm,mathtools}
\usepackage{geometry}
\geometry{a4paper,margin=2.3cm}
\usepackage{microtype}
\usepackage{enumitem}
\usepackage{graphicx}
\usepackage{xcolor}
\usepackage{booktabs}
\usepackage{hyperref}
\hypersetup{colorlinks=true,linkcolor=blue!40!black,urlcolor=blue!50!black}

\setcounter{tocdepth}{2}
\numberwithin{equation}{section}

\theoremstyle{plain}
\newtheorem{teorema}{Teorema}[section]
\newtheorem{proposicao}[teorema]{Proposição}
\newtheorem{lema}[teorema]{Lema}
\newtheorem{corolario}[teorema]{Corolário}

\theoremstyle{definition}
\newtheorem{definicao}[teorema]{Definição}
\newtheorem{exemplo}[teorema]{Exemplo}
\newtheorem{exercicio}{Exercício}[section]

\theoremstyle{remark}
\newtheorem{observacao}[teorema]{Observação}
\newtheorem{contraexemplo}[teorema]{Contraexemplo}
\newtheorem{intuicao}[teorema]{Intuição}

\newcommand{\R}{\mathbb{R}}
\newcommand{\N}{\mathbb{N}}
\newcommand{\Z}{\mathbb{Z}}
\newcommand{\Q}{\mathbb{Q}}
\DeclareRobustCommand{\eps}{\varepsilon}
\DeclareRobustCommand{\vet}[1]{\vec{#1}}
\DeclareRobustCommand{\dif}{\mathrm{d}}
\DeclareRobustCommand{\abs}[1]{\left\lvert #1 \right\rvert}
\DeclareRobustCommand{\norm}[1]{\left\lVert #1 \right\rVert}
\DeclareMathOperator{\grad}{grad}
\DeclareMathOperator{\Div}{div}
\DeclareMathOperator{\rot}{rot}

\newcommand{\base}[1]{\par\medskip\noindent\textbf{Base para relembrar.} #1\par\medskip}
\newcommand{\cuidado}[1]{\par\medskip\noindent\textbf{Quando a hipótese falha.} #1\par\medskip}

\newcommand{\figura}[2]{%
  \begin{figure}[htbp]
  \centering
  \IfFileExists{#1}{%
    \includegraphics[width=0.84\textwidth]{#1}%
  }{%
    \fbox{\parbox{0.8\textwidth}{\centering Figura ainda não gerada:\\ \texttt{#1}}}%
  }
  \caption{#2}
  \end{figure}%
}

\newenvironment{esqueleto}{%
  \par\medskip\noindent\textbf{Esqueleto da demonstração.}\par\nopagebreak
  \begin{enumerate}[leftmargin=1.4em]
}{\end{enumerate}\medskip}

\title{Caderno --- Desigualdade triangular e Cauchy--Schwarz}
\author{}
\date{}
\begin{document}
\maketitle
\tableofcontents

\section*{Notação}

Para $x\in\R$, $\abs{x}=x$ se $x\ge 0$ e $\abs{x}=-x$ se $x<0$. Em $\R^n$, $n\ge 1$, o produto escalar e a norma euclidiana são
\[
\vet{u}\cdot\vet{v}=\sum_{i=1}^n u_i v_i,\qquad
\norm{\vet{u}}=\sqrt{\vet{u}\cdot\vet{u}}.
\]

\section*{Enunciados}

Escreva cada demonstração por completo, com quantificadores e contas. Não há roteiro além do enunciado.

\begin{teorema}[Desigualdade triangular em $\R$]
Para todos $x,y\in\R$, $\abs{x+y}\le\abs{x}+\abs{y}$.
\end{teorema}

\vspace{7.5cm}

\begin{corolario}[Diferença dos módulos]
Para todos $x,y\in\R$, $\abs{\,\abs{x}-\abs{y}\,}\le\abs{x-y}$.
\end{corolario}

\vspace{7.5cm}

\begin{teorema}[Cauchy--Schwarz em $\R^n$]
Para todos $\vet{u},\vet{v}\in\R^n$,
\[
\abs{\vet{u}\cdot\vet{v}}\le\norm{\vet{u}}\,\norm{\vet{v}}.
\]
\end{teorema}

\vspace{9cm}

\begin{teorema}[Desigualdade triangular da norma euclidiana]
Para todos $\vet{u},\vet{v}\in\R^n$,
\[
\norm{\vet{u}+\vet{v}}\le\norm{\vet{u}}+\norm{\vet{v}}.
\]
\end{teorema}

\vspace{8cm}

\end{document}
